% AN \accent READS THE ASSIGNMENTS IN FRONT OF ITS CHARACTER ITSELF.
%
% Once \accent has its character number it goes on reading, and it takes
% blanks, \relax and any assignment before the character it will go over.
% That reading is the accent's own, not the main loop's, so
% \tracingcommands draws no line for any of it -- HSTeX drew one for
% every assignment and every \relax.
%
% Anything that is not one of those ends the wait, and \setbox among them
% is refused where it stands:
%
%   ! Improper \setbox.
%   Sorry, \setbox is not allowed after \halign in a display,
%   or between \accent and an accented character.
%
% which is the same rule, and the same words, as the assignments a
% display alignment reads before its closing $$.
%
% trip line 322 writes `\accent\space"42 \def\^^M{\ } ; \char'101'.
\catcode`\{=1 \catcode`\}=2
\font\ff=cmr10 \ff
\tracingonline=1 \hsize=100pt \hbadness=10000 \hfuzz=100pt
\message{[A assignments between an accent and its character]}
\tracingcommands=2
\setbox0=\hbox{\accent`A \count5=1 \def\z{}\relax B}
\tracingcommands=0
\showthe\count5
\message{[B a non-assignment stops the run]}
\tracingcommands=2
\setbox0=\hbox{\accent`A \count5=2 \kern3pt B}
\tracingcommands=0
\message{[C setbox is refused there]}
\tracingcommands=2
\setbox0=\hbox{\accent`A \setbox1=\hbox{} B}
\tracingcommands=0
\message{[done]}
\end
